\documentclass[fontset=fandol, a4paper, 12pt, openright]{ctexbook} 

\usepackage{geometry}
\usepackage{hyperref}
\usepackage{fancyhdr}
\usepackage{titlesec}
\usepackage{enumitem}
\usepackage{setspace}
\usepackage{xcolor}
\usepackage{amsmath}   % 记得加上这个，因为你正文里有容量计算公式
\usepackage{tabularx} % 记得加上这个，因为你正文里有表格
\usepackage{booktabs} % 记得加上这个，用于表格的三线表样式

\usepackage{tikz}
\usetikzlibrary{positioning} % 提供更好的定位功能支持

% --- 2. 页面与行距 (保持原版宽松布局) ---
\geometry{left=2.5cm, right=2.5cm, top=3cm, bottom=3cm}
\linespread{1.5} 

% --- 3. 元数据与链接设置 (已更新为本书信息) ---
\hypersetup{
    colorlinks=true,
    linkcolor=black, 
    filecolor=magenta,      
    urlcolor=blue,
    % 电子书信息
    pdftitle={力量训练原理},
    pdfauthor={迈克·伊瑟泰尔, 詹姆斯·霍夫曼, 查德·韦斯利·史密斯},
    pdfsubject={用七大核心原理科学提升力量},
    pdfkeywords={力量训练},
    pdflang={zh-CN}
}

% --- 4. 页眉页脚 ---
\pagestyle{fancy}
\fancyhf{}
\fancyhead[CE]{\leftmark}
\fancyhead[CO]{\rightmark}
\fancyfoot[C]{\thepage}

% --- 5. 标题格式 (配合12pt的标准大小) ---
\ctexset{
    chapter = {
        format = \centering\Huge\bfseries,
        number = \chinese{chapter},
        name = {第, 章},
        beforeskip = 10pt,
        afterskip = 30pt
    },
    section = {
        format = \Large\bfseries,
        beforeskip = 20pt,
        afterskip = 15pt
    },
    subsection = {
        format = \large\bfseries, 
        beforeskip = 10pt,
        afterskip = 10pt
    }
}

% --- 7. 图表标题样式 ---
\usepackage{caption}
\captionsetup{
    labelfont={sf},  % 标签使用：无衬线(sf)
    textfont={sf},      % 标题文字使用：无衬线(sf)
    labelsep=space,    % 分隔符使用句点
    justification=centering % 居中对齐
}

\usepackage{enumitem}
\setlist{nosep} % 或者使用 nolistsep，彻底取消列表项之间及列表前后的额外间距

% --- 6. 引用环境 ---
\renewenvironment{quote}
  {\list{}{\rightmargin\leftmargin}%
   \item\relax\small\itshape\color{black}}
  {\endlist}

% --- 书籍信息 ---
\title{\Huge \textbf{力量训练原理}\\\Large 用七大核心原理科学提升力量}
\author{\textbf{【美】迈克·伊瑟泰尔, 詹姆斯·霍夫曼, 查德·韦斯利·史密斯}}
\date{}

\begin{document}

\frontmatter

\maketitle

\chapter*{作者介绍}
\addcontentsline{toc}{chapter}{作者介绍}
\markboth{作者介绍}{作者介绍} % 设置页眉

\section*{迈克·伊瑟泰尔博士}
迈克是美国宾夕法尼亚州费城天普大学运动科学系的教授，曾为中央密苏里大学的教授，负责教授运动生理学、个人训练以及体运能高阶课程三门课程。迈克拥有运动生理学博士学位，担任过美国得克萨斯州约翰逊城奥林匹克训练中心的运动营养顾问。迈克博士作为教练，负责教学员饮食与力量训练，他的学员包括力量举、举重、健美运动员和运动爱好者。迈克曾是力量举、健美和巴西柔术运动员。当大家都在痴迷于有装备力量举比赛时，迈克博士已经多次打破了地区、国家乃至国际无装备力量举纪录。

\section*{詹姆斯·霍夫曼博士}
詹姆斯是天普大学运动科学系的教授，他在东田纳西州立大学获得了运动生理学博士学位，其导师为迈克·斯通博士，研究课题为“推雪撬车对提高橄榄球运动员运动表现的作用”。詹姆斯博士曾作为东田纳西州立大学助理教练、首席运动学家、首席体适能教练和大学举重房管理员，指导了很多橄榄球运动员。他本人也是运动员，在英式橄榄球、美式橄榄球和摔中成绩优秀。

\section*{查德·韦斯利·史密斯博士}
查德·韦斯利·史密斯是 JTS（红坦克训练体系）的创始人及所有者，是在役力量类运动员中成就最高的人之一起初，查德为田径运动员，曾经赢得了两次全美大学生锦标赛铅球项目的冠军。作为力量举运动员，查德绑膝无装备力量举总成绩为1055千克（2325 磅），排名世界第十；套膝总成绩为1010千克（2226磅）排名世界第六。查德在2012 年获得“壮汉职业卡”，也是美国国家橄榄球联盟、美国职业棒球大联盟、终级格斗冠军赛、美国格斗赛事与各类项目奥运选手的教练与顾问。

\section*{关于RP}
RP全称为Renaissance Periodization，是一家专攻饮食与训练的咨询公司。RP 的顾问（包括该书作者）负责为各类客户制订饮食与训练计划，指导运动员达到巅峰水平，帮助职场人员更有精力和改善形体，保持身体健康。自创立公司起，RP 的CEO尼克·肖（NickShaw）便以为客户提供最好的饮食和训练指导为目标。通过聘请数位运动成绩顶尖，又拥有运动学、营养学或生物学博士学位的运动员，尼克组建了一支在健身行业无可匹敌的专家队伍。除了指导训练与饮食，RP也就饮食、训练、周期科学以及各种与身体组成和运动相关的主题持续发表、制作文章与视频等。

\section*{关于JTS}
JTS全名为Juggernaut Training Systems，创立于2009年，其创立之初便以“为所有运动员提供最好的训练”为目标。之后JTS逐渐演变成健身教育中心，各类顶级运动员、教练与研究员聚集于此，为客户和观众提供高质量的内容，包括力量、体适能、营养和机动性训练等。JTS致力为客户提供真正来自于专家的信息内容，帮助运动员实现目标。JTS总部位于加州的尼古湖，是南加州顶级的力量训练中心，有巨大的体育馆，可以进行力量举、举重、体能和“壮汉”等项目，这里同时也是范式运动表现疗法的发源地。 

\chapter*{前言}
\addcontentsline{toc}{chapter}{前言}
\markboth{前言}{前言} % 设置页眉

相信你也一直在搜寻关于如何进行力量训练的内容。我们都曾经试着在网络上搜索文章，但大多数时候总是失望。现在确实有几本关于力量举计划的书籍，而且如果你试过书里的计划会发现：这些书确实值得一看，但这些计划背后的理论来自哪里？身体是如何应对训练的？又如何预测一个计划的效果是好还是坏？到底是什么原理支配着这一切？其实，人体不是由细胞随机组合而成的……人类身体的构建方式大同小异，对训练的反应也几近相同。因此我们设想，关于力量训练，是否能找到一系列的规则或原理？能否将其应用到力量举中？

其实前人已经找到了，而且几乎是适用于所有运动训练的原理。伊萨林（Issurin）、沃邦（Bompa）、哈夫（Haff）、斯通（Stone）、桑兹（Sands）和福赫萧斯基（Verkoshansky）等人已经对运动训练和适应性背后的原理进行了研究。在他们的研究与实践过程中，专项性、超负荷、变式、阶段增益等与训练计划有关的原理已经得到了明确的解释，并催生出了“周期”这个概念。以原理为基础，有逻辑地规划训练，能够帮助运动员获得更大的收益。

自1970年起，世界上几乎每一个奥林匹克训练中心都在使用周期计划，这使得从举重、游泳到体操的训练者的运动表现都发生了跨越式进步。

本书是我们第一次集中讨论力量训练原理对一般运动和力量举的影响。基本原理一共有7条，会在本书中详细介绍，它们对制订力量训练计划来说都举足轻重。

本书能够帮你掌握以下能力。
\begin{enumerate}
\item 深人理解训练，解答疑问，激发学习的欲望。
\item 评价训练计划，根据自身情况调整计划来提高效果。
\item 根据他人对训练的理解，挑选更合适的教练。
\item 根据自己的情况设计基础力量／力量举计划。
\item 根据自己的目标制订计划，收放自如，难简随心。
\end{enumerate}

应用训练原理不仅仅要考虑原理本身，而且还要考虑原理的优先级。7条原理对于训练的重要性是不同的，制订计划时优先参考更重要的原理，不仅可以提高训练效果，而且可以避免将时间与精力浪费在细枝末节上。

本书将讲解力量训练计划／周期设计的7条基础原理，尤其还会关注原理在力量举中的应用。根据重要性，我们会将原理进行排序，从而帮你调整训练计划。因此，无论你的计划是精心设计还是较为简略的，都可以提高训练效果。

该书将会深挖这7条原理，探索它们背后的科学，寻找实际应用它们的方法同时解答应用过程中出现的常见问题。不是每位读者都可以看懂这本书，希望你们能将训练看成一个过程，勤动脑筋；不要满足于“怎么做”，而是要去探寻“为什么”这么做，也不要仅仅满足于机械地执行教练的指令。读书学习是一个漫长的过程，只有这样，你们才会发现这本书是多么有趣。

该书不会出现以下事情。
\begin{enumerate}
    \item “能不能简单点？”追求进步的过程不会过分复杂，但也不会简单到一目了然，力量／力量举周期训练需要不断思考，没有人可以直接从概念推算出结果。
    \item “能不能少努力点？”如果你想尽可能不努力，又想获得更多的收益，那你不用看这本书了。通过训练来让自已变强是一件严肃的事情，每次训练都需要你认真对待。你可以调整计划，降低对训练的投人，这由你自已取舍………·但本书所有的建议都基于你想获得最大收益，并且你愿意为之付出相应的努力。
    \item “能不能你说我做？”这本书不会提供许多制订好的计划，也没有精确的组次和重量百分比。原理很直观，但需要你深入思考，并根据自身的情况来调整实践，寻找最大收益。这本书更像一个工具，帮助你学习原理、理解训练，而不是仅仅只是教会方法。所有的原理都源于研究周期训练的最新资料与文献。
\end{enumerate}

希望你们能够真正地去理解力量训练的“为什么”，为什么有些计划有效果，而有些是徒劳？为什么有些计划对一些人有效，对另一些人却无效？尽管内容专业，但我们会以轻松易懂的方式来讲解，借助专业术语，由浅人深地剖析书中内容。为了方便技术讨论，本书的第一章将会介绍一些术语。

话不多说，我们开始吧！


\tableofcontents

\mainmatter

% \setcounter{chapter}{4}
\chapter{重要术语}
在我们真正开始深入学习力量举训练理论之前，读者朋友们最好和我们一同熟悉或复习一下基本术语，这对接下来的学习至关重要。你在之前的训练和学习中可能已经了解了一部分，但还存在一些让人疑惑的术语，我们今天就来把这些障碍扫除吧。

\textbf{强度（Intensity）：} 
又称“绝对强度”，指的是训练时需要对抗的阻力，通常用磅（1磅 $\approx$ 0.4536 千克）或千克表示。例如，负重200千克深蹲一次的强度要高于195千克10次，即使后者使得训练者完全力竭。

\textbf{相对强度（Relative Intensity）：} 
任何一组训练中，向心阶段最接近肌肉力竭的程度。假设你负重200千克可以做6次，这组训练只做了3次；负重180千克可以做9次，这组训练做了8次，那么后者的相对强度更高。

\textbf{容量（Volume）：} 
一次推举、一组动作、每次或每周训练的总机械功。从技术层面来讲，容量的计算方法：
\begin{equation*}
    \text{容量} = \text{总受力} \times \text{总行程}
\end{equation*}
精确计算容量的公式：
\begin{equation*}
    \text{容量} = \text{组数} \times \text{次数} \times \text{重量} \times \text{杠铃行程}
\end{equation*}
但多数情况下，容量是让个体训练者比较单个动作的概念。在这种情况下，行程可以从公式中去掉。因此，容量计算公式简化为：
\begin{equation*}
    \text{容量} \approx \text{组数} \times \text{次数} \times \text{重量}
\end{equation*}

\textbf{频率（Frequency）：} 
一段时间内的训练次数，通常指一周内训练几次。

\textbf{动作选择（Exercise Selection）：} 
每次训练计划内所选择的动作。比如我们说练腿，但更精确的描述可能是深蹲和直腿硬拉。

\textbf{训练单元（Training Session）：} 
单次训练内所有动作的总和，单位为“个”。每个训练单元通常包括热身、正式训练和收尾，可能还会包括冷身。

\textbf{轻训单元（Light Session）：} 
指故意设计成没有超负荷（详见第四章）的训练，完成难度相对较低。其目的是促进恢复与适应，同时缓解技术退步、肌肉损失和力量降低。常用手段是降低容量或强度。

\textbf{休息日（Rest Day）：} 
没有训练的日子。

\textbf{小循环（Microcycle）：} 
重复一个训练单元中的所有动作类型的时间单位，通常是指1周。示例安排如下：

\begin{table}[h]
\centering
\caption{典型小循环安排示例}
\begin{tabularx}{\textwidth}{XXXXXXX}
\toprule
周一 & 周二 & 周三 & 周四 & 周五 & 周六 & 周日 \\
\midrule
卧推 & 深蹲 & 休息日 & 过顶推举 & 硬拉 & 休息日 & 休息日 \\
\small{上半身辅助} & \small{下半身辅助} & & \small{上半身辅助} & \small{下半身辅助} & & \\
\bottomrule
\end{tabularx}
\end{table}

\textbf{累进期（Accumulation Phase）：} 
由一系列连续的小循环组成，训练难度会逐渐增加，包括强度增加、容量增加或二者一起增加。

\textbf{减载期（Deload）：} 
指一个全部都由轻训组成的小循环，目的是减轻疲劳，保持已达到的训练适应。

\textbf{中循环（Mesocycle）：} 
由一系列小循环组成，目的是为了实现特定的训练适应。典型的中循环由累进期和减载期组成，通常持续1个月左右。

\textbf{进减比（Accumulation-to-Deload Ratio）：} 
指的是在一个中循环内，累进期与减载期所用时间的比例。最典型的进减比是 3:1。

\textbf{训练模块（Training Block）：} 
一个训练模块通常包含1$\sim$3个中循环，其中每个中循环的训练目标相似或相同。

\textbf{大循环（Macrocycle）：} 
由一系列训练模块按照特定顺序组成。一个力量举的大循环通常由增肌模块、增力模块、冲刺模块、比赛日组成。

\textbf{体重管理：} 
\begin{itemize}
    \item \textbf{增重：} 提高肌肉量，高容量训练配合热量盈余，持续1$\sim$3个月。
    \item \textbf{维持：} 保持体重。常用语增力期和冲刺期，应至少持续1个月。
    \item \textbf{减重：} 降低脂肪量，高容量训练配合热量赤字，持续1$\sim$3个月。
\end{itemize}

\textbf{力量（Strength）：} 
生成力的最大能力，实验室以“牛顿”为单位，健身房常用“千克”或“磅”。

\textbf{专项力量（Specific Strength）：} 
在特定动作中能够输出极限力量的能力。通常将专项中最薄弱的环节视为专项力量。

\textbf{周期（Periodization）：} 
按照特定逻辑顺序安排的训练变量组合，目的是让训练者适应训练、降低受伤概率和提高特定时间的竞技状态。

\textbf{机动性（Mobility）：} 
综合运用技术、力量和活动度的能力。常见不足情况包括：
\begin{enumerate}
    \item 力量和活动度足够，但技术不足。
    \item 力量和技术足够，但活动度不足。
    \item 活动度和技术足够，但力量不足。
\end{enumerate}

\chapter{训练原理简介}

力量/力量举训练不是想到什么就练什么，然后祈祷老天爷把你变强壮。每一个训练计划的制订都应参考一定的训练原理，然后通过训练来验证这些原理是否有效，能否提高你的运动成绩。实践是检验真理的唯一标准，我们需要探索身体对训练的反应，然后根据得到的结果来安排训练计划。只有这样，我们才能更进一步。

假设外星人天生躺着休息就可以变强壮，活动一下就会影响他们的进步，每一次训练都会损害他们的运动表现，那么，这个星球的力量训练者的最好训练计划就是在不训练时尽可能多地休息。但如果地球上的人类只休息而不训练的话，身体只会越来越弱。人类训练者与那个外星球的训练者相反，需要训练来打破人体稳态提供刺激，产生新的生理适应，从而提高力量。

因此，对人类力量训练来说，“超负荷”原理至关重要。如果不持续挑战训练者的生理极限，提高训练的难度，训练者的力量就不可能增大。也许宇宙里存在另一个星球，那个星球的人的恢复能力远比我们低，即使是最基础的训练也需要两周才能恢复，那么他们的最佳训练频率可能就是每三周训练一次。但这样的训练计划能在我们身上发挥作用吗？当然不能，因为人体有特定的生理机制，对训练和休息会以相应的方式进行反应。所以用于指导人类力量训练的计划必须根据其生理机制来制订。

人类力量训练有 7 条训练原理。依照这 7 条原理，我们能够精心制订计划，提高训练收益。这些原理都不是偶然发现的、不是编的、不是发明创造的，而是前人经实证研究验证，并 100\% 基于训练者的生理反应和运动表现的真实数据得来的。忽视这几条原理可能会导致训练适得其反，或是无法预测运动员的成绩。

以下 7 条训练原理按重要性由上到下排列，并包含它们各自的子原理（主原理中的细节要点）：

\begin{enumerate}
    \item \textbf{专项性 (Specificity)}
    \begin{itemize}
        \item 子原理一：训练模式的兼容性
        \item 子原理二：定向适应
    \end{itemize}
    \item \textbf{超负荷}
    \item \textbf{疲劳管理}
    \item \textbf{刺激—恢复—适应 (Stimulus-Recovery-Adaptation, SRA)}
    \item \textbf{变式}
    \item \textbf{阶段增益}
    \item \textbf{个体差异}
\end{enumerate}

这些原理基于人体生理系统得出，应用于训练过程，所以在计划中应用这些原理可以提高寻训练者的运动表现和成绩。但这 7 条原理对训练者成绩的影响是不同的，排名越靠前，影响越大。单是基于“专项性”和“超负荷”原理的计划效果也会很好，但恰当应用的原理越多，计划的效果则越好，顶级的训练计划会体现出所有的原理。

下图展示了各原理的重要性排序，越靠近底部、图形面积越大的原理越重要。“专向性”涵盖了整张图，意味着这是最重要的、必须最先应用的应用原理，在力量训练中应当有最高优先级，我们会在第三章详细解释原因。每一条训练原理都会单独用一章来详细讲解。

\begin{figure}[htbp]
    \centering
    \begin{tikzpicture}[x=1cm, y=0.7cm, font=\small]
        % --- 1. 左侧刻度与垂直标题 ---
        % 刻度线与百分比
        \draw[dotted] (-2.5, 0) -- (-1, 0) node[left=1.5cm] {0\%};
        \draw[dotted] (-2.5, 10) -- (-1, 10) node[left=1.5cm] {100\%};
        
        % 垂直排列的标题文本
        \node[rotate=0, align=center, text width=1em] at (-2, 5) {力\\量\\训\\练\\原\\理\\优\\先\\级};

        % --- 2. 绘制分层叠块 (灰度方案) ---
        % 注意：每层的高度和灰度根据原图比例做了微调
        \fill[gray!10] (0, 0.2)   rectangle (5, 2.5); % 专项性 (最底层)
        \fill[gray!60] (0, 2.7)   rectangle (5, 5.0); % 超负荷
        \fill[gray!40] (0, 5.2)   rectangle (5, 6.8); % 疲劳管理
        \fill[gray!15] (0, 7.0)   rectangle (5, 8.0); % 刺激—恢复—适应
        \fill[gray!25] (0, 8.2)   rectangle (5, 9.0); % 变式
        \fill[gray!55] (0, 9.2)   rectangle (5, 9.7); % 阶段增益
        \fill[gray!35] (0, 9.9)   rectangle (5, 10.3); % 个体差异 (最顶层)

        % --- 3. 右侧标签与虚线指向 ---
        % 坐标位置根据叠块的中点计算
        \draw[dotted] (5.2, 1.35)  -- (6.5, 1.35) node[right] {专项性};
        \draw[dotted] (5.2, 3.85)  -- (6.5, 3.85) node[right] {超负荷};
        \draw[dotted] (5.2, 6.00)  -- (6.5, 6.00) node[right] {疲劳管理};
        \draw[dotted] (5.2, 7.50)  -- (6.5, 7.50) node[right] {刺激—恢复—适应};
        \draw[dotted] (5.2, 8.60)  -- (6.5, 8.60) node[right] {变式};
        \draw[dotted] (5.2, 9.45)  -- (6.5, 9.45) node[right] {阶段增益};
        \draw[dotted] (5.2, 10.1)  -- (6.5, 10.1) node[right] {个体差异};
    \end{tikzpicture}
    \caption{训练原理优先级叠块图}
\end{figure}

下面介绍原理的每一章的行文结构大致如下。

\begin{enumerate}
    \item 该原理在运动科学领域的定义。
    \item 该原理在力量举领域的定义及拓展。
    \item 解释该原理排名的原因，并与其他原理进行比较。
    \item 讲解和示范如何合理运用该原理。
    \item 实例讲解该原理应用不足的后果。
    \item 实例讲解该原理过度应用的后果。
    \item 总结该原理与力量举的关系。
    \item 参考资料与拓展阅读。
\end{enumerate}

\chapter{专项性}

\section{科学定义}

专项性原理是力量训练里最重要的原理，其科学定义也是最直接明确的。确切来说，专项性指训练应当让负责该项运动任务的基础系统产生负荷，从而刺激该系统。也就是说，专项训练是训练负责完成目标动作的系统，可以是训练需要募集的肌肉，可以是训练需要调动的神经系统，也可以是训练动作本身。专项性的定义很广，我们会在下文中详细讨论。

\section{在力量举领域的定义}

在力量举领域，专项性的定义是“针对能提高力量举成绩的系统所进行的训练”，其范围从广到狭而有所变化。广义的专项性是指任何能够增加发力肌肉体积的训练，狭义的专项性是指能够提高这些发力肌肉基础力量的训练，而最具体的定义就是能够改善实际比赛中动作执行与提高运动成绩的训练。

肌肉体积这个大分支可以再分成不同的小分支，比如力量举分支或健美分支，但它们各自仍然具有力量举或健美共有的专项性。沿此路径再往下分，就可能是专门练上胸或练下胸的具体训练实践。但是耐力却是另一个分支，虽然它和肌肉体积来源相同——身体的总体健康状况。耐力和健美、力量举并不在同一个分支上。如果你进行的耐力训练越多，你就离力量举和健美越远。

% ==========================================
% TODO: 专项性谱树 图片占位符
% 等您提供截图后，可以使用 tikz 绘制或直接用 \includegraphics 插入
% ==========================================
\begin{figure}[htbp]
    \centering
    \begin{tikzpicture}[x=1cm, y=1cm, font=\sffamily] % 使用无衬线体(黑体)使图表更现代
        % --- 1. 马拉松 (最左) ---
        \fill[gray!55] (-5.5, 6) -- (-3.8, 6) -- (-2.5, 0.5) -- (-3.1, 0.5) -- cycle;
        \node[above, text=black!80] at (-4.65, 6.2) {马拉松};

        % --- 2. 800米跑 (左二) ---
        \fill[gray!85] (-3.5, 6) -- (-1.5, 6) -- (-0.9, 0.5) -- (-2.2, 0.5) -- cycle;
        \node[above, text=black!80] at (-2.5, 6.2) {800米跑};

        % --- 3. 100米跑 (正中) ---
        \fill[gray!30] (-1.2, 6) -- (1.2, 6) -- (0.6, 0.5) -- (-0.6, 0.5) -- cycle;
        \node[above, text=black!80] at (0, 6.2) {100米跑};

        % --- 4. 健美 (右二) ---
        \fill[gray!85] (1.5, 6) -- (3.5, 6) -- (2.2, 0.5) -- (0.9, 0.5) -- cycle;
        \node[above, text=black!80] at (2.5, 6.2) {健美};

        % --- 5. 力量举 (最右) ---
        \fill[gray!55] (3.8, 6) -- (5.5, 6) -- (3.1, 0.5) -- (2.5, 0.5) -- cycle;
        \node[above, text=black!80] at (4.65, 6.2) {力量举};

        % --- 中间覆盖的白色标签框 ---
        % fill=white 实现了覆盖切断背景的效果
        \node[fill=white, inner xsep=15pt, inner ysep=8pt, text=black!80, font=\large\bfseries]
            at (-3.55, 3) {耐\quad 力};
        \node[fill=white, inner xsep=10pt, inner ysep=8pt, text=black!80, font=\large\bfseries]
            at (3.55, 3) {肌肉体积};

        % --- 底部总体标签 ---
        \node[font=\LARGE\bfseries, text=black!80] at (0, -0.4) {总体健康状况};
    \end{tikzpicture}
    \caption{专项性谱树}
    \label{fig:speciality-tree}
\end{figure}

从\autoref{fig:speciality-tree}可以看出，力量举和健美训练者都可以从增肌训练中受益，因为肌肉体积这个分支最终同时指向了这两种运动。肌肉体积这个概念范围的宽窄决定了增肌训练是否仅是力量举这一项运动的专项训练。专项性是一个范围的概念（可以把它理解成一张树状图），所以会有各种各样的分支训练，有一些训练对力量举的专项性较低，只能对提高力量举成绩起一点点作用，或者根本就没用。每周做好几次柔韧性训练不会影响力量举成绩，但对力量举训练也没太大好处，只是沾点边罢了。有很多的训练模式与力量举专项无关，有些甚至会起反作用，最好的例子就是耐力训练。

我们可以按专项性原理来假设一个系统，将力量举的训练分成 4 类：
\begin{enumerate}
    \item 能够根据专项特点，针对性地提高力量举成绩的训练（大重量三项训练）。
    \item 大体上能够提高力量举成绩的训练（增肌与增力训练）。
    \item 与提高力量举成绩沾边的训练（高阶活动度训练）。
    \item 无关甚至是损害力量举成绩的训练（耐力训练）。
\end{enumerate}

这告诉我们，专项性的定义范围会从广泛到狭窄，甚至还会出现负向的专项性训练。举例来说，狭义的专项性训练有卧推技术训练，广义的专项性训练有胸肌增肌训练，负向的专项性训练有会导致力量下降的超量耐力训练。这引出了专项性两个子原理中的第一条——训练模式的兼容性。

\section{子原理一：训练模式的兼容性}

专项性的第一个子原理是训练模式的兼容性。在对力量举训练的研究中，人们发现有的训练形式能极大地提高成绩，有的效果没那么好，有的根本没作用，有的甚至会起反作用。根据训练的实际作用来设计训练计划对提高训练者的成绩至关重要。除此之外，有些训练模式之间更能够相互兼容。某种模式的训练可以帮助训练者达成多种训练目标。比如，通过举重训练让运动员增加的力量也可以让他在游泳时更有力量。在兼容性的另一端，有些训练模式虽然能改进运动员某方面的能力，但也会损害某些特定的运动表现，比如让马拉松训练者增肌太多就会降低他的力量与体重比，最终导致他跑得更慢。

训练模式的兼容性对运动项目的影响是很有趣的话题，可拓展的地方很多，但详细分析每种运动将超出本书的讨论范围。接下来，我们会集中讨论力量举中最常见的训练模式之间的兼容性。

力量训练于其他训练模式之间的兼容性由低到高排列如下：
\begin{itemize}
    \item 力量训练与耐力训练
    \item 力量训练与活动度训练
    \item 力量训练与爆发力训练
    \item 力量训练与增肌训练
    \item 基础力量训练与专项力量训练
\end{itemize}

我们会依次讨论这些训练模式，探索它们之间的兼容性与冲突，从而帮助我们更好地制订计划。

\subsection{力量训练与耐力训练}

除低强度有氧运动之外，在多数情况下，耐力训练会极大地阻碍力量训练。为提高耐力而进行的刻苦训练会干扰许多与力量举相关的生理系统：

\textbf{肌肉流失：} 耐力训练会分解肌肉组织，所以训练者如果长期进行高强度的耐力训练会极大地增加肌肉流失的风险。肌肉量下降会导致能够产生力量的收缩性组织减少，继而影响训练者的力量输出。

\textbf{疲劳：} 所有的训练都会积累疲劳。只积累疲劳却无法提高运动表现的训练是有害的，疲劳会影响训练者的恢复和之后的运动状态。如果你在练力量举的同时还在进行健美训练，尽管这不是最高效的方法，但至少健美练出来的大肌肉能够帮你输出更多的力量，这在一定程度上也能提高你的力量举运动表现。但是耐力训练既不能让你增肌，也不能提高你的力量，但会让你积累大量的疲劳。耐力训练对于力量训练几乎没有任何益处。即使你的体型实在太差了，想通过运动来改善，耐力性运动也不应是你的首选。耐力训练通常是高容量训练，积累的疲劳会非常多。所以力量举训练者即使选择下棋，也比做耐力训练好，因为下棋至少不会积累疲劳。

\textbf{肌纤维类型改变：} 人体骨骼肌纤维包含快肌纤维和慢肌纤维。快肌纤维对力量举更重要，因为接收神经信号之后，快肌纤维的出力速度更快，而且每个单位横截面积的力量更大。不仅如此，与慢肌纤维相比，一般人肌肉中快肌纤维的含量更高，在大重量训练的刺激下，增长速度更快。因此，在身体条件允许的情况下，训练者尽可能提高体内快肌纤维的含量，从而能够让他在更短的时间内提高力量举成绩。如果坚持这么做，随着时间的推移，你会变得越来越强壮。

另一方面，与快肌纤维相比，慢肌纤维在收到神经信号后，收缩速度更慢，单位横截面积的力量更小，一般人体内的慢肌纤维含量也更低。此外，慢肌纤维增粗、变肥大的潜力也更小。所以，从力量举的角度来说，提高快肌纤维含量是更优的选择。这就引出了我们要讨论的问题，与其他类型的训练相比，耐力训练使肌肉纤维的快、慢肌比例改变的能力是最强的。耐力训练不会改变肌肉纤维的类型，但介于两种肌纤维类型之间的中间纤维（有些甚至会同时具有两种肌纤维的特性）会表现得更加像慢肌纤维。这种类型特征的改变确实是可能发生的，但无论是由快变慢还是由慢变快，都要花数月的时间才会发生。从现实角度来讲，耐力训练会降低人体内快肌纤维的储备，从而影响力量举训练者的运动表现。

\textbf{神经系统改变：} 运动单位指的是运动神经（由脊柱发出，负责支配肌肉）与其所支配的肌肉细胞（纤维）的合称。肌纤维性质的改变会使得整个运动单位发生改变，也就是说运动神经也会为了肌纤维性质改变之后的活动而做出调整和优化。耐力训练不仅会影响局部运动神经，其范围还会延伸到脊柱，甚至是大脑。这种改变会让人更加擅长低强度、持续性的活动，更加不擅长高强度、爆发性的活动。耐力训练导致的大脑和神经系统的改变会深刻影响力量举的运动表现。

因为会导致疲劳积累、肌肉功能和神经功能的改变，耐力训练对力量举来说基本上是百害而无一利，因此力量举训练者要尽可能减少耐力训练，这样才能使自己力量举训练的效果最大化。当然，如果你有自己的理由，或者你还从事其他运动，那么在做出权衡之后，当然也可以安排耐力训练。但这是一种权衡之后的选择。耐力训练可以和其他一些运动同时进行，但最好不要与举重或力量举一起进行，这对力量举专项训练没有好处。

\subsection{力量训练与活动度训练}

活动度是指单个或多个关节的活动范围。举个例子，肱二头肌弯举时肘关节的活动范围是单关节活动度，完成硬拉动作所涉及到的多个关节的活动范围就是多关节活动度。活动度有 4 个基本因素：

\begin{itemize}
    \item \textbf{肌肉活动度：} 肌腹长度、肌肉在安全范围内的延展性以及最大拉伸长度。肌肉活动度受到肌肉组织本身以及神经系统的限制。
    \item \textbf{结缔组织活动度：} 筋膜和连接肌腹、肌腱、韧带的结缔组织在安全范围内的延展性。
    \item \textbf{关节活动度：} 关节运动时所通过的运动弧或转动的角度。关节结构的位置和形状（比如股骨头的形状、膝盖半月板的厚度）在某种情况下决定了该关节的活动度。
    \item \textbf{邻近组织的影响：} 在某些情况下，肢体会碰触到身体邻近的其他部位继而限制了其活动度。例如，无论你做深蹲时能蹲多低，当你的腘绳肌和小腿碰到一起时，就没法继续往下了，这种干扰对肌肉发达的训练者来说尤为明显。
\end{itemize}

力量举中所需要的活动度可以通过以下三种方式来进行训练：
\begin{enumerate}
    \item 进行力量举常规关节活动度（Range of Motion, ROM）的训练，比如深蹲、卧推、硬拉等。
    \item 进行力量举超常规关节活动度的训练，比如“高杠早安”、水牛杆卧推、超程硬拉等。
    \item 进行专门用于增加活动度的拉伸训练。
\end{enumerate}

前两种方法对力量举训练没有负面影响，但有明确证据表明拉伸训练（尤其是经常性、高强度的训练）会积累疲劳，干扰力量适应。如果你的活动度很差，使你在训练过程中产生了安全的隐患，或导致你无法完全发挥实力，而且上面提到的前两种方法已经无法提供有效的帮助了，那么我们可以妥协一下，额外安排一些专门用于增加活动度的拉伸训练，这对整体训练效果还是有益的。但是如果你的活动度变得超过完成力量举常规动作行程所需要的程度，反而有可能干扰你的最佳适应速度。

超出需求的活动度训练还会导致另外两个问题：第一，极端的活动度会影响关节安全，有研究表明，超灵活以及欠缺灵活度的关节受伤概率更高。不仅如此，极端的活动度会降低在动作底部时的组织紧绷程度（比如深蹲到底部的反弹），从而影响伸展—收缩循环（Stretch Shortening Cycle, SSC）的效果，这就在一定程度上降低了训练者的力量表现。第二，有足够的证据表明活动度训练会暂时降低力量和爆发力的输出（短期的影响，大概持续几分钟到几个小时）。所以，如果活动度训练结束之后紧接着进行力量举训练，那么力量举训练一定会大受影响。因此，我们建议在完成力量训练后再进行活动度训练，或者两种训练的间隔时间延长一些。

因此，我们尝试性地总结出以下的活动度训练方案：
\begin{itemize}
    \item 保质保量地完成常规的力量举训练动作。
    \item 定期进行超常规关节活动度的力量举训练动作的训练，尤其是当你感觉活动度有问题时。
    \item 如果前两种方法不够，就额外安排一些增加活动度的拉伸训练，从而提高力量举专项所需的活动度，但拉伸训练与力量训练之间的间隔时间要长一些。
    \item 保证必须的基础活动度就行，力量并不会因为你的活动度特别好而额外增加。
\end{itemize}

\subsection{力量训练与爆发力训练}

力量指的是输出力的能力。爆发力是指在短时间做出最大功的能力。力量训练与爆发力训练的兼容性比上面两种情况好，但我们仍然能在其中找到一些问题。

爆发力训练也需要时间和精力，我们可以将这部分时间和精力放在更有益于力量举的事情上来，比如安排更多的力量训练、恢复时间和增肌训练等。爆发力本身对力量举的帮助也不大，在大重量的力量举训练中，动作完成速度是很慢的，爆发力对其基本没有影响，力量输出才是决定性因素。增强完成动作时的意志力对力量举训练很重要（比如挺过粘滞点），这时使用的负荷与力量训练中所使用的负荷通常是一致的。但爆发力专项训练使用的重量通常在 30\% 1RM 左右，这比力量举比赛中使用的重量强度要低很多，所以对力量举的专项性也比较低。

爆发力输出受疲劳的影响非常大，所以在相对低容量的训练时更容易输出爆发力。即使是轻微的训练容量的增加也会导致肌纤维的类型向慢肌转化，同时还会造成神经疲劳，从而降低爆发力输出。但是力量训练需要中等容量（明显要比爆发力训练的容量高），增肌时更是要使用高容量训练。因此爆发力训练所采取的训练容量对提升力量的效果并不好，长期进行爆发力训练还会影响增肌，这更是与力量举训练不兼容。有趣的是力量的英文是 Strength，Power 是爆发力的意思，但“力量举”的英文却是 Powerlifting，而不是 Strengthlifting。

力量训练与爆发力训练确实会相互影响，如果你只想提高力量举成绩，那就需要舍弃一部分爆发力训练，把更多的精力放在增肌和增力训练上。我们会推荐某些项目的运动员同时训练力量与爆发力，但是力量举不属于这些项目。

\subsection{力量训练与增肌训练}

增加肌肉量可以增强肌肉的收缩能力，继而直接增强力量输出，但是增肌训练仍然会影响（微小的程度，但仍值得注意）力量举成绩，其主要体现在以下方面：

通常来说，增肌效果好的训练容量会很高，这会导致惊人的疲劳积累，让肌纤维的类型发生转变，干扰神经系统的功能，最终让训练者疲于应对超负荷训练，影响其力量举成绩。增肌训练对力量举训练者的帮助很大，但效果的显现有一定的延迟性。增肌训练结束之后，训练者通常会降低容量，在低疲劳的情况下进行增力训练，此时增肌训练的优点才会显现出来。

过度训练与力量举专项无关的肌肉会干扰力量举的训练。肱二头肌、内侧三角肌和小腿的训练对力量举很有帮助，但过度训练这些部位的肌肉会积累大量疲劳，反而得不偿失。肩膀外翻看起来确实很帅，这对健美运动很重要，但却基本无法提高力量训练者的力量水平。过度训练与力量举无关的肌肉，不仅不会有助于力量举，而且还会积累疲劳，浪费宝贵的训练和恢复时间。

力量举训练者当然可以力型兼备，而且主要肌群之间保持平衡可以预防损伤。但在增力和冲刺阶段进行增肌训练就没什么作用了，甚至还会影响训练者的最终成绩。补充一点，如果你的训练周期拉得很长，在增力和冲刺阶段进行增肌训练的坏处就显得没那么大了，但这么做依然可能会影响你的最佳表现。

\subsection{基础力量训练与专项力量训练}

基础力量训练与专项力量训练是目前来说互相干扰最小的两种训练模式。基础力量指的是使用身体肌肉组织来发力的基础能力，而专项力量指的是在单次特定动作中，比如深蹲、卧推、硬拉，身体输出力量的能力。

一个能很好说明两者之间区别的例子就是超大重量哑铃卧推与竞赛卧推。有些训练者更擅长哑铃卧推，或者器械卧推、窄推、地板推等，但这些都不能以准确的算法来转化为训练者 1RM 竞赛卧推时的表现；或者有些训练者做 5RM 卧推时能使用的重量非常大，但却不擅长 1RM 卧推的重量。

通常，这里还涉及一些技术问题。在极限重量下，好的技术给训练者带来的帮助非常大，而不当的技术会严重阻碍训练者的运动表现。更深入一点来讲，接近你负重极限的重量（1RM）会让你的技术趋近于崩溃（无法保持技术）。所以，有的运动员即便基础力量水平很高，但如果技术很差，也无法在专项动作中发挥出良好的专项力量。

还要补充一点，训练者在低容量、低疲劳的状态下才能输出自己最大的专项力量，而力量训练最好采用中等容量。因此，想要同时表现出最佳的基础力量和专项力量，基本是不大可能的，但这之间的差距并不大。

训练模式具有兼容性的意思不是说所有以 1RM 为目标的训练都是好的，也不是说所有不以 1RM 为目标的训练都是不好的，这与专项性有关。专项训练确实能够提高专项成绩，但是其他更广泛的训练模式也可以提高专项成绩，比如基础力量与增肌训练，只要运动员产生了这样的需求，我们就可以在合适的时间采用这样的训练模式。但是有一些专项性不高的训练模式（比如活动度训练）对力量举既没好处也没坏处，就需要你特殊地对待、谨慎地进行。还有一些专项性较高的训练（比如耐力训练），如果与你正在进行的训练兼容性不高，就需要尽可能地避免。

\section{子原理二：定向适应}

专项性的第二个子原理是定向适应。这个概念在关于力量举训练的讨论中很少被提到，但其实它对制订力量举训练计划很重要。

定向适应的定义在很大程度上依赖于专项性的定义。具体来说，专项性是指如果你想要提高一个系统的能力，你就需要对这个系统或其基础系统施加压力。而定向适应的意思是我们所施加的这些压力需要按照特定的顺序来安排。作为专项性的子原理，定向适应需要我们按顺序、持续性地来施加压力，提供专项刺激，这样才能最大限度地提高运动表现。

在这里，顺序指的是训练刺激要以特定的次序来安排，并且要有时间的限定，而不是随机安排训练。举个例子，如果你想提高每组重复 5 次的深蹲能力，我建议你最好每周都安排几次每组深蹲 5 次的特定训练，而不是每 3 周甚至 6 周才做一下每组重复 5 次的深蹲训练。

定向适应与生理学息息相关。首先，人体系统的适应能力是有限的，尤其是经过严格训练的人体系统。其次，虽然每一次训练都能带来微小的适应，但在训练刺激停止后，这些训练适应能持续的时间却取决于训练安排的顺序和强度。所以，人体系统的适应能力有上限（必要的时候需要放弃某项适应来达成另一个专项性更高的适应），以特定顺序安排且强度合理的训练得到的训练适应保持的时间最长、效果最好。

举个例子，如果你的目标是完成每组 5 次的深蹲，你需要达成一系列的适应，包括神经、肌肉和结缔组织等方面，从而帮助你实现目标。根据专项性原理，直接安排每组 5 次的深蹲训练是比较好的方法。但我们假设你不仅想提高每组 5 次的深蹲能力，还想提高每组重复蹲 15 次的能力，这时每组深蹲训练所安排的动作次数就应该在这两个次数之间。然而，每组 5$\sim$15 次的深蹲训练对于想要提高 5 次 1 组深蹲的运动表现来说，肯定不如全部都安排每组 5 次深蹲训练的效果好。本来只需要向着每组 5 次来适应的人体，由于定向适应的第一个局限（人体有限的适应能力），有了另外一个适应方向——每组深蹲 15 次。为了向每组深蹲 15 次的适应转变，一些人体系统，比如肌纤维和/或神经系统等，势必会产生相应的变化，这就会影响每组深蹲 5 次的适应。除此之外，由于定向适应的第二个局限（训练安排顺序），如果你每一次都安排每组 15 次的深蹲训练，那么你之前每组 5 次的深蹲训练所带来的适应就会倒退，并朝着每组 15 次的适应方向转化。训练刺激因此而缺少持续性，导致训练者无法巩固训练效果，每组 5 次的训练适应在日后无法得到很好的维持。

以上只是举了一个简单的例子，实际的力量举训练要复杂得多。我们见过不少这样的训练计划，它们在一周内同时安排了每组重复 5 次的训练、大重量复合组训练和每组重复 20 多次的递减组训练。其实每一种方法都有各自的好处，但是因为定向适应的关系，如果简单地把它们全部加起来使用，这样做的结果就是运动员的身体产生了各种各样的适应，但是每一种都不突出，也就是俗话说的“博而不专”。不仅如此，在综合训练阶段（进行每组重复 3 次的训练来促进神经适应、进行每组重复 20 次的训练来促进增肌适应）获得的训练适应会更容易消退。一旦我们想将目光集中在一个更明确的目标时（比如提高每组重复 5 次的运动表现），就只剩下坏处了。

在后文里，我们会讨论阶段增益，那时我们会提出更好的方法来利用定向适应原理，从而帮助训练者制订更精密的计划。在增肌阶段，我们会安排一系列特定的训练，并且使整个计划的时间足够长（1$\sim$N个月），从而尽可能地维持训练适应。在基础增力阶段，我们会安排与增肌阶段不同的训练刺激，利用增肌期获得的肌肉，强化训练者的力量表现。在比赛之前，也就是冲刺期，我们会巩固训练者在增力期获得的力量，让训练者做好准备，迎接最大的负重试举。我们使用多种训练方法来提高训练者的力量举成绩。这些训练方法经过精密安排，能够同时提高训练者多方面的素质，并且不违反定向适应原理。

\section{训练原理重要性排序}

我们根据每条原理对力量举训练的重要性和对整体计划安排的影响程度来给 7 条训练原理排序。其中专项性排第一，这是毫无争议的，因为它优先于其他的训练原理。假设一个计划以超负荷为主，那么训练者可能需要每周跑几百千米；一个计划如果以阶段增益为主，那么它可能会从基础有氧阶段快速过渡到比赛阶段；一个计划如果以疲劳管理为主，那么它可能天天都是进行轻复合跑步和单车训练。但是，这些训练计划不会让你的力量举产生一点点的进步！因为这些是“铁人三项”的专项计划。所以我们说专项性优先于其他的训练原理，只有在保证专项性的基础上，帮助训练者提高与某项运动相关的各方面的素质，才能提高其在某项运动中的成绩。力量举也是如此。我们会在本章后面讨论专项性应用不足的例子。在训练计划中，专项性的应用不足是非常严重的错误，影响着计划设计和训练者进步。

\begin{figure}[htbp]
    \centering
    \begin{tikzpicture}[x=1cm, y=0.7cm, font=\small]
        % --- 1. 左侧刻度与垂直标题 ---
        \draw[dotted] (-2.5, 0) -- (-1, 0) node[left=1.5cm] {0\%};
        \draw[dotted] (-2.5, 10) -- (-1, 10) node[left=1.5cm] {100\%};
        \node[rotate=0, align=center, text width=1em] at (-2, 5) {力\\量\\训\\练\\原\\理\\优\\先\\级};

        % --- 2. 绘制分层叠块 (灰度方案) ---
        \fill[gray!10] (0, 0.2)   rectangle (5, 2.5); % 专项性 (最底层)
        \fill[gray!60] (0, 2.7)   rectangle (5, 5.0); % 超负荷
        \fill[gray!40] (0, 5.2)   rectangle (5, 6.8); % 疲劳管理
        \fill[gray!15] (0, 7.0)   rectangle (5, 8.0); % 刺激—恢复—适应
        \fill[gray!25] (0, 8.2)   rectangle (5, 9.0); % 变式
        \fill[gray!55] (0, 9.2)   rectangle (5, 9.7); % 阶段增益
        \fill[gray!35] (0, 9.9)   rectangle (5, 10.3); % 个体差异 (最顶层)

        % --- 3. 右侧标签与虚线指向 ---
        \draw[dotted] (5.2, 1.35)  -- (6.5, 1.35) node[right] {专项性};
        \draw[dotted] (5.2, 3.85)  -- (6.5, 3.85) node[right] {超负荷};
        \draw[dotted] (5.2, 6.00)  -- (6.5, 6.00) node[right] {疲劳管理};
        \draw[dotted] (5.2, 7.50)  -- (6.5, 7.50) node[right] {刺激—恢复—适应};
        \draw[dotted] (5.2, 8.60)  -- (6.5, 8.60) node[right] {变式};
        \draw[dotted] (5.2, 9.45)  -- (6.5, 9.45) node[right] {阶段增益};
        \draw[dotted] (5.2, 10.1)  -- (6.5, 10.1) node[right] {个体差异};
    \end{tikzpicture}
    \caption{训练原理重要性排序：专项性涵盖了整张图的最底部}
\end{figure}

\section{专项性原理的正确应用}

\subsection{比赛前的动作选择}

由于变式原理的好处（我们会在第七章详细讨论），在一个训练大循环内，不是每一个中循环都需要安排竞赛动作。但是，随着比赛日期临近，安排的训练动作必须要越来越接近竞赛动作。这样做有 2 个原因，我们会在接下来的内容中一一解释。但因为变式原理，我们不会在一个大循环内的每个中循环都安排竞赛动作，所以“每次都要练竞赛动作”在我们这里也不是最佳方案。我们将在第七章讲解变式的时候，详细讨论每次都练竞赛动作的优缺点。

不是每次训练都要安排竞赛动作，但是随着比赛日期临近，安排的训练动作必须要越来越接近竞赛动作，以下是 2 条主要原因：

\textbf{1. 肌肉与结缔组织适应}

在试举极限重量的时候，杠铃的重量会由人体精密地分布在各个肌肉上，同时还有肌腱、骨头和韧带参与承重。在刚开始一个新大循环训练的时候，我们会设定一系列的目标（如用前蹲、腿举来强化股四头肌，改善深蹲的力线），但在比赛日时，测试的却是训练者的力量输出能力和结缔组织承重的能力。通过训练，人体会以一种精密的方式发展自身的承重能力。研究发现，经过数周的训练，即使训练者的肌肉量没有增加，其肌腱和肌肉的超微结构（肌腱中胶原蛋白排列和肌肉羽状角）也会根据专项动作而改变。变式训练能够使训练者的力量输出能力和人体承重能力在一个较长的时间范围内得到优化，但这种能力的提升通常会出现在训练数月之后，而如果想在短期（数周）内强化这两种能力，直接训练竞赛动作是最好的方法。因此，随着比赛日期的临近，尤其是在冲刺期（赛前最后一个中周期），我们要把训练动作逐渐调整为竞赛动作，从而优化人体适应。

随着比赛日期的临近，我们不仅要把训练的主要内容逐渐调整为比赛动作，辅助动作也应该越来越具有专项性。但为了保证辅助动作不过分干扰专项动作，我们需要在赛前调整辅助动作，从而能够更好地在短期内（冲刺期）将训练适应转化为比赛成绩。举个例子，在非赛季，“碎颅者”和过顶臂屈伸这两个动作锻炼肱三头肌的效果都很好；但是随着比赛日期的临近，我们更推荐窄距卧推和肩推，因为这两个动作与竞赛动作的角度更相似，募集的相同肌肉更多。赛前的辅助动作越接近比赛动作，在短期内就越容易转化为比赛成绩。

\textbf{2. 神经适应}

肌肉收缩需要神经发出信号。竞赛动作和辅助动作对于神经系统的不同影响可以从 3 个方面来讲：
\begin{itemize}
    \item \textbf{运动单位：} 在局部区域，对于竞赛动作和辅助动作来说，神经系统的工作方式大致相同，但仍有一些区别。举个例子，无论是卧推中的肱三头肌收缩，还是“碎颅者”中的肱三头肌收缩，每个被激活的运动单位的工作方式都是完全一样的。但是，两个训练动作在肱三头肌上激活的运动单位不是完全一样的。卧推时，肱三头肌上有一部分运动单位没有激活，但在做“碎颅者”时被激活了。反之亦然。
    \item \textbf{肌肉：} 竞赛动作和辅助动作对于神经系统的不同影响不仅体现在运动单位上，在肌肉方面也有所区别。比如胸肌对卧推的作用很大，对“碎颅者”的作用却一般。
    \item \textbf{动作：} 不同的动作会募集不同的肌群，并且募集肌肉的时机、发力时长和收缩力量也都不相同，这会导致在不同动作下，神经系统的活动不同。窄距卧推和满握竞赛卧推的动作模式基本一致，而且募集的运动单位和肌肉发力模式也比较类似，但是运动单位的募集顺序和肌肉激活度也会有所不同。甚至即使是训练同一个动作，如果训练者使用了不同的技术，神经系统的活动差异也会很大。赛前我们不仅要调整训练者运动单位和肌群的神经功能，还要调整大脑功能，从而提高训练者的成绩。
\end{itemize}

力量举作为专项运动，在一定程度上确实需要追求极致的技术。不管你在备赛早期练了何种动作，为了在比赛中发挥出最佳实力，一定需要安排与规定比赛动作相应的技术训练。

现在我们从现实的角度来看看，如果赛前几周不练竞赛动作会造成什么后果：
\begin{enumerate}
    \item 结缔组织受伤概率增高，因为它们不适应竞赛专项动作的力线（尤其是比赛中冲刺最大力量的时候）。
    \item 力量输出能力较差，直接参加比赛会导致身体无法完全利用增肌期获得的肌肉，神经系统与肌肉的羽状角也未对竞赛动作进行相适应的调整和优化。
    \item 试举时容易出现技术失误（比如硬拉时臀先起），因为赛前没有进行相应的技术训练。
\end{enumerate}

我们说的这些后果都是前人在实践经历中所遭受过的。在 21 世纪早期，有许多训练者（多数来自西部杠铃）有一种“骄傲”，他们会在赛前什么动作都练，却唯独不练硬拉，在比赛中的硬拉热身是他们自备赛期以来的第一次硬拉。硬拉带来的疲劳感确实是非常高的（我们会在第五章详细剖析），但通过完全不做硬拉训练来规避疲劳积累并不是最佳的选择。总之，简单来讲，随着比赛日的到来，你的训练需要越来越接近竞赛动作。

\subsection{大循环下的辅助动作}

专项性在竞赛动作的训练中会非常明显的体现，但在涉及没有包含竞赛动作的训练时，我们也应该有意识地将专项性合理地运用到其中。在利用专项性来安排辅助动作时，我们通常会遇到两方面的限制：

\textbf{1. 大局观}

无论你做什么样的辅助练习，最终目的都是为了提高你的力量举成绩。除了活动度训练、技术训练、伤病管理（这确实也算是一种辅助手段，但通常不把它当作一种直接的体育训练），其他所有力量举辅助训练的目标，除了提高力量举成绩之外，都只能聚焦在两个方面：增加肌肉量和强化动作力量。

\textbf{肌肉量：} 当你做增肌训练时，请确保你训练的肌肉真的是需要被强化的。如果你的目标是强化三角肌中束，而辅助动作却安排了大量的小腿训练，那就违反了专项性原理，因为这些辅助训练对实现你目标的作用并不大。刺激、恢复、适应的资源是有限的，训练某些部位的肌肉就意味着你会错过另外一些部位的训练，或者影响到它们的恢复。你应该去训练那些对竞赛动作最有效的肌肉，而不是那些对你的目标帮助并不大的肌肉。通常，对这些肌肉所做的辅助训练没有什么特别助益。有些关于辅助训练的传闻是毫无信服力的。例如，三角肌后束确实对卧推有所帮助，但与胸肌的作用相比呢？所以胸肌的优先级一定高于三角肌后束。

\textbf{动作力量：} 就像肌肉量一样，动作力量的训练也需要制定优先级。如果你在硬拉时，从地面启动的速度非常快，但在动作过程中逐渐开始驼背，那你就需要提高背部力量，而不是增加深蹲训练。寻找自己的弱点，然后选择合适的动作，这对高阶训练者来说尤为重要，我们会在后文详细解答这个问题。但弱势部位训练太多也可能产生问题，因为恢复速度会跟不上。我们会在第五章详细解析最大可恢复容量的概念，教你选择合适的训练容量。无论如何，辅助动作不应该随便安排，它需要更直接的有助于提高训练者的动作力量或者与这些训练动作相关的肌肉。这种避免随意安排辅助训练的需要就引出了辅助训练的第二个限制因素：避免非定向训练。

\textbf{2. 避免非定向训练}

所有不能促进专项适应的训练都是一种浪费，原因是与其做无用功，不如休息，或者安排更具专项性的训练。除了一些特殊的考量外，如果辅助训练无法增加力量举相关部位的肌肉量或强化相关动作的力量，那么这些训练可能就是没用的甚至是一种浪费，以下这些便是例子：
\begin{itemize}
    \item 小腿训练、高次数腹肌训练、肱二头肌训练。
    \item 过度的背部训练（有时可能是有用的）。
    \item 各种类型的有氧锻炼。
    \item 各种瑜伽、普拉提和活动度训练。
\end{itemize}
简单来说，如果你对自己做这个训练的原因还不明确，这就不是一个好兆头！

\subsection{学习与训练适当的技术}

除了一些荒谬的力量比赛，力量举的评判标准是非常清晰明确的。专业的力量举协会在比赛判罚的标准上会保持高度统一。我们的日常训练最终需要转化为比赛成绩，所以我们的训练需要和竞赛动作保持一致或接近。

助力带可以帮助提高硬拉的重量、组数和次数，暂停式卧推也许可以强化某些训练者的肩部力量。诸如此类的训练调整在非赛季时使用是可以的，但随着比赛日期临近，尤其是赛前一个月的训练，我们的训练动作必须要和竞赛动作保持完全一致。赛前一个月开始使用规范的竞赛动作来训练是比较恰当的时机，赛前一周才开始就太晚了，因为那个时候如果训练强度太高，就很难在比赛中较好地应用竞赛技术，所以在冲刺期一开始就切换为竞赛动作是比较合适的。

有一个经典的例子可以用来说明，太晚开始竞赛动作的训练会带来许多困难。一位训练者按照教练的要求，在比赛前的两三周一直进行弹力卧推训练，这是教练基于训练者弹力卧推的 1RM 所制定的每组 3 次的训练。但在比赛的前几天，教练却突然要求训练者停止弹力卧推，改为和竞赛动作一致的训练，即杠铃触胸停稳再推起。训练者一下子乱了阵脚，并且很可能无法完成训练，因为竞赛标准卧推可比弹力卧推难多了。这样错误的训练计划会打乱训练者的训练节奏，影响训练者的心态。如果在赛前才突然将训练动作调整为竞赛动作，不仅会降低训练者的训练重量，而且执行两种动作所需要的技术也有区别。从弹力卧推改为竞赛标准卧推、从砸地反弹硬拉改为竞赛标准硬拉，都需要花几周的时间来仔细地调整技术，训练者在这个过程中会觉得动作陌生、别扭。而在赛前冲刺期，训练者的训练重量非常大，这时候最需要避免的就是让训练者觉得动作陌生和不适应。

总结一下，我们需要学习适当的技术，日复一日地训练它，特别是在赛前的几个月，这样才能在比赛中发挥出最佳水平。

\section{专项性原理应用不足}

专项性原理是力量举训练中最重要的原理，为了提高训练的效率，我们需要将专项性原理恰当地应用在计划中，避免过度应用或者应用不足。在本书中，我们会专门讨论专项性的应用不足和过度应用，以及导致的不良后果。接下来我们首先讨论制订力量举计划时 3 个常见的专项性应用不足的例子。

\subsection{使用与力量举不兼容的训练模式}

由于训练模式的兼容性和定向适应这两个原因，并且身体从训练中恢复和适应的能力也是有限的，所以，所有的外部训练（指非力量举专项训练）都有可能干扰力量举训练，影响运动员的进步。各种外部训练形式对力量举训练的干扰是不同的，有些较高，有些较低。

判断外部训练对力量举干扰程度的 4 个主要指标如下：
\begin{enumerate}
    \item 外部训练的容量。
    \item 外部训练的影响/阻碍。
    \item 外部训练的运动模式。
    \item 外部训练的时机。
\end{enumerate}

外部训练的容量越高，就越会对训练者进行力量举训练时的能量水平和训练后的恢复造成影响（继而降低训练成果）。除了力量举训练外，如果你每周还打 2 次网球，每次 1 个小时，这其实对你的力量举训练没造成什么干扰。但如果你每周打 6 次网球，每次都持续 2 个小时，这就很有可能对力量举训练有害了。

外部训练对训练者生理系统的干扰程度也是重要的指标。如果除了力量举训练之外，训练者每周还额外进行 6 个小时温和的瑜伽练习，这只会对你的生理系统增加一些压力，并不会让你的身体出现磕碰、瘀血，或者伤病。但是如果你还喜欢综合格斗（MMA）或者巴西柔术（BJJ），那么它们的训练强度、对人体造成的干扰，以及在运动中可能造成的伤害，你都需要慎重考虑。

外部训练的运动模式是非常重要的指标，此前我们在讨论训练模式的兼容性（作为专项性的子原理）时，就已经详细讨论了这个问题。与健美相比，马拉松训练对力量举训练的干扰更大，因为从生物化学、肌肉和神经层面来说，马拉松训练对力量举训练的阻碍更大。

尽管较少被关注，但外部训练的时机也会影响它对力量举训练的干扰程度。我们从两个角度来讨论训练时机：在训练日/训练周内，外部训练与力量举训练的接近程度，以及在一个力量举大循环内，外部训练的参与时机。为了减少外部训练的干扰，将它与力量举训练之间的时间间隔安排得越长越好，这也有助于提高每次力量举训练的效果，并且促进训练者的恢复和适应。举个例子，训练者早上踢足球，然后到晚上再进行力量举训练，要比 2 个项目紧接着练要好。以周为单位来算的话，比较好的安排是在周六、日练综合格斗、巴西柔术等，周二至周五练力量举，这样当你进行力量举训练时，身体已经能较好地从外部训练中恢复过来。长期来看，在一个大周期内，越接近比赛日，训练计划中的专项性要越强。也就是说，如果你必须要做外部训练，但又想力量举训练的效果最大化，我们建议在大周期的前期多做外部训练，而随着力量举计划的推进，不断降低外部训练的占比，将目标更加集中在力量举的训练、恢复和适应上。

苏联的一些力量训练者如果被选入举重项目的奥林匹克国家队，他们就需要接受这么一条明确规定：禁止参与任何可能会影响举重成绩的身体活动。教练们常用的做法就是在一年的大多数时间内，将训练者带到集训营训练，让他们没有机会去参与任何外部训练与活动。但如今，多数人只是将力量举作为爱好，而不是职业，所以让训练者严格避免所有的外部活动是不太可能的。其实，也没有这个必要，只要这些外部活动能够充实生活，增加乐趣，就是值得的。但是，根据自己想要在力量举这项运动中取得的成绩，爱好者们可以根据以上的 4 个指标来决定选择何种外部训练，决定外部训练的活动量和参与时机，衡量力量举训练与外部训练的效果，从而更有效地提高力量举成绩。训练者不应该迷迷糊糊地参与外部训练，浪费自己的精力，或者在罪恶感中进行外部训练。

\subsection{使用非专项性基础体能储备训练}

许多知名力量举训练体系都详细解释了基础体能储备（General Physical Preparedness, GPP）的作用，并大力推崇它。基础体能储备主要是帮助训练者保持足够好的体型，从而能够在接下来的专项训练中获得最大化的效果。

而对于基础力量举训练，也被称为专项体能训练（Special Physical Preparedness, SPP），业界基本上已经达成了共识。许多 GPP 的训练方法都非常流行，但是也存在一些 GPP 训练在一定程度上违反了专项性原理。GPP 训练的整体目标是提高运动员专项训练的效果。对于力量举来说，GPP 训练就是提高训练者的肌肉量，为提高力量打下基础，并且提高训练者对 SPP 训练的恢复和适应能力。以当代周期化理论的视角来看，GPP 训练属于增肌期，在此阶段，力量训练者会提高目标肌群的训练容量和每组训练重复的次数（通常是每组重复 6$\sim$10 次），为增力期（SPP 训练阶段）做好肌肉量基础与训练能力的准备。力量举训练的专项性，以及想打下牢固训练基础等原因，训练者在 GPP 训练阶段应选择专项性非常高的动作，可以只改变竞赛动作的一些变量（如选择前蹲、窄距卧推、超程硬拉等），这些动作的训练对训练者后期成绩的提升有很高的迁移性。

如果力量举的 GPP 训练目的是增肌和提高专项能力，那么将 GPP 训练的内容安排为壶铃摇摆或推雪橇训练对力量举又有什么用呢？我们对此的回答是直截了当的，没太大作用。壶铃训练对力量举所需的肌群和动作的专项性都不够，同时此项训练又很难实现超负荷，对力量举动作的迁移性也不高。推雪橇训练通常是用来提高足球和橄榄球运动员某些运动能力的，比如加速能力等，但这是进行足球和橄榄球运动所需要的体能储备，而不是力量举。虽然动作不同，刚开始练也很有趣，但是像壶铃、高强度间歇性训练（HIIT）和推雪橇这些训练模式在力量举的 GPP 训练里没有太高的地位，因为它们对力量举的专项性不高，有很多更具专项性、更高效的动作可以替代它们。不是说力量举很特殊，才排斥这些训练模式的，而是因为每项运动都有其自身的专项性，GPP 训练一直都需要与专项性相结合。还记得我们在本书一开始的树状图吗？GPP 是在树枝上，而不是树干上。

\subsection{设计计划时没有使用或遵循专项性}

专项性应用不足的最后一个例子就是训练者选择了很多训练模式。这些训练模式虽然有时候看起来有用，但实际上极有可能与力量举背道而驰。在制订训练计划时，我们需要挑选训练方法、训练模式和训练动作，这其中的每一样都必须要有理有据，而不是随随便便凭着自己的感觉选择。我们至少需要照顾到以下的一个方面：
\begin{itemize}
    \item 肌肉量
    \item 力量
    \item 冲刺
    \item 恢复
    \item 适应
    \item 伤病预防
    \item 技术
\end{itemize}

而且这些方面的训练最终都需要与提高力量举成绩挂钩。比如，肌肉量指的是要提高力量举训练时需要使用的肌肉，而不是每一块肌肉都要练，技术指的是提高力量举技术，而不是其他运动的技术。

举个例子，机动性训练可以帮助训练者更好地掌握训练技术，从而让训练者在进行全程动作训练时更加安全和舒适。但是机动性训练做得太多也会带来问题，额外的机动性训练不会对力量举有任何好处。如果你的目标是练好力量举，那么你计划里的每一项安排都必须有理有据并与力量举直接相关，而不是九转十八弯之后才和力量举沾一点边。简单来说就是，如果你一时无法找到现在的训练和组次安排的依据，那么多数情况下，你现在的训练对力量举的专项性就不会高，我们建议你就不要再使用这个计划了，计划必须是与力量举的专项性密切相关的。

\section{专项性原理过度应用}

在上面我们谈到了专项性的应用不足，但其实专项性也可以被过度应用。专项性原理与变式原理在很多方面是相对应的，所以专项性的过度应用可能就是变式的应用不足，反之亦然。我们会在第七章介绍变式原理时详细讨论这个问题，但我们先来介绍几个专项性过度应用的例子。

\subsection{常年使用“竞赛重量”训练}

常年使用“竞赛重量”训练指的是训练者在一整年的训练中每个月、每周都大量安排每组重复 1$\sim$3 次的接近极限重量的训练。这个训练模式很有趣，因为它违反了许多基本的训练原理，但这里我们只讲 4 个原因，解释为什么这种训练违反了专项性原理。

\textbf{1. 训练过期}

即使一个训练模式有效，但如果长时间只应用同一种训练模式也会降低训练者的适应效率，我们将这种现象称为过期。也就是说，如果你在一段长时间内都使用同一种训练模式，你的训练效率会越来越低，继而变成一种精力的浪费。这就引到了变式原理上，我们会在第七章介绍变式原理的时候详细讲解训练过期的生理原因。因此在某些极端情况下，专项性和变式是互相矛盾的，专项性过度应用也可以称为变式应用不足。

训练变式太少会降低人体对训练刺激的反应，继而降低训练效果。变式可以通过多种形式来实现，包括不同的动作安排、组间休息时间、训练重量和训练组次。例如，仅仅是改变训练组次，其他的安排保持不变，就能让训练进步一直持续下去。因此，长期使用竞赛重量（每组重复 1$\sim$3 次的大重量）训练的效果会越来越低，但此时只要将你的组次安排改变一下就又可以继续获得进步了。

训练过期并不是特指长时间使用大重量（每组 1$\sim$3 次）训练，只要你的计划都一直使用同一个固定的组次，也会有这个问题。大重量计划之后立刻接上小重量训练周期就可以有效缓解训练过期的问题。也就是说，相比于其他方法，大小重量组合的训练模式是比较有效的变式，能够帮助规避训练过期的问题。

\textbf{2. 错过增肌训练}

力量举训练有一个很重要的部分，那就是训练者不仅需要提高现有肌肉的力量，还需要增肌。除非你刚开始练力量举时就又矮又壮，而且在比赛中你会遇到的对手都是那种又高又瘦的新手，你比他们天生高了几个体重级（至少 5 个），否则在一般情况下，你都需要调整身体肌肉量的比例，从而让自已变得更有竞争力。通常来说，确实是肌肉量越大，力量的增长空间就越大。

而一直使用我们刚才说的“竞赛重量”训练就会错过增肌期，因为增肌需要的是高容量训练，而非大重量训练。在容量相同的情况下，确实是训练重量越大，增肌效果越好，但是超过 80\% 1RM 的训练重量会让训练者非常疲劳，你很难在这个强度下坚持训练下去。所以，如果训练者只进行大重量训练，容量就上不去，也就无法保持快速进步。这种训练方式会浪费你作为力量举训练者的潜力，长久如此，你的竞争对手就很有可能把你远远地甩在后面。

那能不能在一次训练中先练大重量，然后再降低重量来提高训练容量，从而同时获得两种训练的好处呢？在一定程度上这是可以的，但前提是你的训练计划要对大重量和高容量有优先级的区分，一个为主，另一个为辅。如果你的训练安排是 50\% 的时间进行大重量训练，50\% 的时间进行高容量训练，那么这两种训练结合起来的效果勉强还行，但是绝对谈不上出色。因为这样的训练结合会降低训练的增肌效果，五五开之中的高容量训练会在两方面影响大重量训练。首先，高容量训练会导致单次训练就大量积累疲劳，训练者很难快速恢复，那么在下一次训练中就很难再使用同样的大重量，同时还会影响到接下来的容量训练。其次，采用大重量训练会让训练者的神经系统适应单次最大发力的方向，降低了增肌期所需要的多次数训练的神经适应能力。因此，从专项性原理来看，高容量训练与大重量训练之间是有矛盾的，如果不能处理好两者的关系，最终会导致两种训练都无法产生最好的结果。所以我们还是建议把目光放得长远一点，让每个模块的训练目标都更加明确，我们会在“阶段增益”一章详细讨论这个问题。

\textbf{3. 错过基础力量训练}

长期使用竞赛重量（超大重量）进行训练不仅会影响增肌，它还会造成一个不良后果，并且这个后果是与我们的认知相违背的，那就是这样的训练模式也会影响我们的力量增长。力量增长需要的不仅仅是足够大重量的训练（$\ge$ 75\% 1RM），同时还需要有足够的训练容量（组数 $\times$ 次数 $\times$ 重量），只有这样才能最大化地训练适应，提高训练效果。

采用每组 3$\sim$5 次的训练（使用更接近 80\% 1RM 的重量，而不是高于 90\% 1RM 的重量）更能使训练者保持进步的速度。这样的计划能使训练更可持续，让重量和容量都有所兼顾。重量太大会积累太多疲劳，影响训练效果，从而无法稳固地打下力量基础。每一次无意义的大重量推举都会积累不必要的疲劳，影响你的能量水平和恢复能力。

\textbf{4. 过度强调连续短期的备赛能力}

常年使用极限重量训练，你或许能在力量举比赛前用几周的时间就调整好备赛状态。你只需要比平时的训练稍微降低点辅助动作的容量，并暂停 1$\sim$2 周的大重量训练就可以上场了。似乎你可以通过这种方法来常年保持竞赛状态。但从另一方面来看，细致规划与安排高容量、大重量训练的人却需要数周甚至数月的调整才能进入比赛状态。

这样看来，似乎常年使用竞赛重量进行训练确实也不错。但我相信一定有人会问：在现实的生活中，什么人会只有这么短的备赛时间？或者什么样的狂人会每个月都去打比赛？美国伟大的力量举教练路易·西蒙斯（Louie Simmons）曾经说过，他可以在几周内就让训练者进入比赛状态。但这又如何呢？对于各位读者来说，这世界上有那么多顶级的力量举比赛可以去参加吗？这种连续短期的备赛训练会影响训练者的长足进步。我们的目标是不断突破自己，而不是让自己长期备赛，所以这种连续短期的备赛能力也算不上什么优点。

\subsection{变式动作太少}

就像长期使用固定次数区间的“竞赛重量”训练（每组 1$\sim$3 次的大重量训练）会限制变式原理能带来的训练效果一样，太过限制训练计划内动作的选择也会降低计划的效果。因此，使用不同的变式组合会有至少两方面的好处：一方面是变式自身带来的好处，另一方面是变式可以根据运动员的需求做出调整。

\textbf{1. 变式自身的好处}

过度使用任何一种训练模式都会导致人体的适应停滞，此时我们可以通过引入新的训练模式来解决这个问题。只要新加入的训练模式与其他的训练原理兼容（超负荷、专项性、疲劳管理），这些新动作就可以让训练者继续进步。举例来说，当训练者已经进行了几个中循环的竞赛卧推训练之后，他只要把训练动作换成窄距卧推，就能进一步增肌与增加力量。而当他接着使用了 1$\sim$N 个中循环的窄距卧推之后，再换回竞赛卧推，成绩就又可以更进一步。这样的方法在几个小周期内就可以看到效果，因为神经和肌肉会快速适应新动作，并在之前的基础上更进步。

不仅如此，因为训练者在中间使用了数个中循环的窄距卧推，竞赛卧推也变成了一种新动作，也能够利用变式自身的好处，让训练者持续提高成绩。这种“删除 + 替换”的模式，即每间隔一段时间换一种训练动作的方法，能让训练者的身体对新的动作产生不适应，从而有效地预防训练者的进步停滞，提高训练效果。

\textbf{2. 定向变式}

训练变式是提高训练效果很好的工具，而且我们可以使用定向变式来进一步发挥变式的作用。运动员先天的遗传基因和后天训练经验存在差异，他们的肌肉、动作与技术的优缺点各不相同。有的运动员深蹲能力很强，原因是他的臀肌和竖脊肌异常强壮，但另一位运动员出色的深蹲能力强可能却是因为他强壮的股四头肌。虽然训练竞赛动作确实会对训练者的弱点进行弥补（例如，假设一位训练者的腘绳肌很强壮、股四头肌弱小，那么他练习深蹲动作肯定会对股四头肌造成更大的刺激，股四头肌的适应速度也会更快，最终有可能会追上腘绳肌），但仍然会有例外。

训练者通常会调整技术来规避弱点。例如，以臀肌发力为主导的训练者深蹲时通常会将屁股往后坐一点，这样的做法并不是生物力学上的最佳状态，但却可以最大化他们的深蹲成绩。而他们之所以这样做，其实是一种对自己臀肌强而股四头肌弱的妥协。持续这么练下去，不仅会降低训练者的进步速度，还会让他们的股四头肌与臀肌的差距越来越大。这个问题对高阶训练者来说非常重要，因为他们的弱点和优势非常明确；而对于初阶和中阶训练者来说，情况又不一样了，因为中级训练者还需要发展力量，新手的主要目标在于各方面一起进步。

还有一个常见的问题：人体运动的动力链中有一些部分比另外的部分能够承受更高的训练量。假设某位训练者的股四头肌每周需要 20000 千克的训练容量来增肌或者增力，但他的下背部只能承受每周 15000 千克的训练容量，如果超过这个量就会快速积累疲劳，大幅度增加受伤的概率。所以，也就是说你的深蹲训练也只能承受 15000 千克的训练容量，除非你不要腰了。

想要解决上面提到的问题，动作变式是必不可少的训练方法。举个例子，如果你的深蹲动作是以腘绳肌为主导，那么我们就把你的竞赛深蹲改为高杠深蹲，这样股四头肌的受力就会更多、刺激更大，并且也能保持腘绳肌不发生退步。如果一个训练者的卧推主要依靠胸肌发力，并且他的“锁定状态”（由骨骼关节代替肌肉承受压力的情况）很严重，我们就会建议他多练窄距卧推，以强化肱三头肌，然后临近正式比赛的时候再换回宽距卧推，这样就能比较好地解决训练者的“锁定”问题了。如果你的背部肌肉弱而腿部肌肉强，进行架上拉训练就可以有效提高背部的训练容量，强化背部，而且因为杠铃不是从地上拉起，也很好地规避了腿部的过度训练。

变式有这么多的好处，所以辅助动作训练在力量举里特别受欢迎，按理说也应该如此。但还是有不同的声音存在。有很多影响力大的训练者和教练极度推崇“坚持基础训练”，即在训练计划中不加入任何变式。他们只推荐别人进行比赛三项动作的训练（深蹲、卧推、硬拉）。尽管一些超级精英运动员和高阶运动员在比赛前只练竞赛三项动作，并且也取得了不错的成绩，但我们仍然坚持认为，如果全年都只进行三项动作训练并不是个好方法。

\subsection{辅助动作与竞赛动作太相似}

使用变式动作的目的是提供直接、有效和全新的训练刺激。因此，变式动作必须要有的放矢，并且产生足够的刺激来增肌和增力。新选择的动作要和竞赛动作有足够的区别，这样它们作为变式才能发挥效果。

执行力量举计划时，常见的一个错误是辅助动作的形式与功能都和竞赛动作太相似了，这样，辅助动作的刺激可能就会不足，也就无法完全发挥它作为变式的优点，以实现辅助动作真正的目的。举个例子，我们为了强化股四头肌会安排前蹲，但此时的前蹲并不是蹲起来就可以了，而是要特意去强化股四头肌。如果你用竞赛动作的方式蹲起来，就会变成身体前倾的“早安深蹲”，这样又会变成了以后链肌群为主导的深蹲，失去了前蹲强化股四头肌的作用。身体前倾的“早安深蹲”与竞赛深蹲相比，除了杠铃位置不一样，从生物力学角度上看基本是一模一样，所以它作为变式是无法提供新的刺激的。在练前蹲的时候，最好的方法应该是身体保持直立，尽量用股四头肌发力。在几个中周期的训练之后，你的股四头肌就会得到强化，继而可以帮助提高竞赛深蹲的成绩。简而言之，要保证辅助动作与竞赛动作的不同，这样才能完全挖掘变式动作的优点。这样的训练既能使训练者直接受益（利用前蹲来强化股四头肌），又能给训练者提供新的刺激（与竞赛动作不同的肌肉激活模式，从而推动进步，防止训练停滞）。

\section{总结}

用力量举的行话来说，专项性原理的意思并不是要所有的训练都采取大重量单次的训练，而是说使用的训练最终能够提高力量举专项成绩。换言之，训练要以竞赛动作需要募集的肌肉与肌肉力量为第一优先级。在任何备赛的大循环中，随着比赛日的到来，我们要逐渐以直接的专项性训练（严格竞赛规则下的极限重量试举）为最优先，逐渐降低辅助性训练，比如增肌训练和基础力量训练等。

\subsection*{重点}
\begin{itemize}
    \item 专项性是一个范畴，囊括了能够转化成比赛成绩的一切训练。训练的专项性越强，转化为比赛成绩的效果越好。
    \item 力量举训练的专项性需求包括了增肌训练、增力训练、技术训练、备赛冲刺训练、恢复、适应和伤病管理。
    \item 定向适应是一个过程，我们重复使用训练刺激来强化适应。训练应当有的放矢，而不是在一次训练中什么都练，最后变得博而不专。
    \item 非专项训练或者与力量举进步不直接相关的训练都可以适当舍弃，我们可以用这些时间和训练机会来做更加专项性的训练，或者用来恢复和/或适应。
    \item 专项性不是说全年的训练都要完全模拟比赛情况。力量举训练者不是只能进行极限重量单次/两次的训练，而是通过动作、容量和强度的组合来强化最终的运动表现。
\end{itemize}

\section{参考资料与拓展阅读}

\textbf{专项性的定义与相关内容}
\begin{itemize}
    \item \textit{Principles and Practice of Resistance Training}《抗阻训练的原理与实践》
    \item \textit{Periodization 5th Edition Theory and Methodology of Training}《周期化训练理论与方法论（第五版）》
    \item \textit{Science and Practice of Strength Training}《力量训练科学与实践》
    \item \textit{Essentials of Strength and Conditioning}《力量与体适能训练精要》
    \item \textit{Training Principles: Evaluation of Modes and Methods of Resistance Training - a Coaching Perspective}《训练原理：从教练的角度评估抗阻训练的模式与方法》
    \item \textit{Transfer of Strength and Power Training to Sports Performance}《力量与爆发力训练与运动表现的迁移性》
\end{itemize}

\textbf{训练效果的迁移性}
\begin{itemize}
    \item \textit{Training Transfer: Scientific Background and Insights For Practical Application}《训练的迁移性：科学背景与实际应用》
\end{itemize}

\textbf{骨骼肌肌肥大的专项性}
\begin{itemize}
    \item \textit{Non-uniform Response of Skeletal Muscle to Heavy Resistance Training: Can Bodybuilder Induce Regional Muscle Hypertrophy}《骨骼肌对大重量抗阻训练的非一致性反应：健美运动员可以局部增肌吗？》
    \item \textit{Whole Body Muscle Hypertrophy From Resistance Training: Distribution and Total Mass}《抗阻训练带来全身肌肥大：分布与人体总肉量》
    \item \textit{Non-uniform Muscle Hypertrophy: Its Relation to Muscle Activation in Training Session}《非一致肌肥大与肌肉激活的关系》
    \item \textit{The Adaptations to Strength Training: Morphological and Neurological Contributions to Increased Strength}《力量训练的适应性：增力的形态学与神经学因素》
    \item \textit{Regional Hypertrophy}《局部增肌》
\end{itemize}

\textbf{训练类型的互相干扰}
\begin{itemize}
    \item \textit{Concurrent Training: A Meta-Analysis Examining Interference of Aerobic and Resistance Exercises}《并行训练：用荟萃分析检查有氧训练与抗阻训练的干扰》
    \item \textit{Concurrent Strength and Endurance Training: From Molecules to Man}《力量与耐力并行训练：从细胞到人类》
\end{itemize}

\textbf{疲劳与运动表现}
\begin{itemize}
    \item \textit{Interactive Processes Link the Multiple Symptoms of Fatigue in Sport Competition}《交互过程：体育竞赛中疲劳的多重症状》
    \item \textit{Muscle Fatigue: What, Why and How it Influences Muscle Function}《肌肉疲劳影响肌肉功能的现象、原理和方式》
    \item \textit{Unraveling the Neuro-physiology of Muscle Fatigue}《肌肉疲劳的神经生理学》
    \item \textit{A Comparison of Central Aspects of Fatigue in Sub-maximal and Maximal Voluntary Contractions}《比较极限与次极限肌肉自发收缩导致的疲劳》
    \item \textit{Principles and Practice of Resistance Training}《抗阻训练的原理与实践》
    \item \textit{Periodization 5th Edition Theory and Methodology of Training}《周期化训练理论与方法论（第五版）》
    \item \textit{Tapering and Peaking for Optimal Performance}《获取最佳成绩的冲刺与收尾训练方法》
\end{itemize}

\textbf{活动度训练与运动表现}
\begin{itemize}
    \item \textit{The Effects of Stretching on Performance}《拉伸对运动表现的影响》
    \item \textit{Effect of Acute Static Stretch on Maximal Muscle Performance: A systematic Review}《系统评估急性的静态拉伸对肌肉表现的影响》
    \item \textit{The Effects of Stretching on Strength Performance}《拉伸对力量的影响》
\end{itemize}

\textbf{运动目的的专项适应}
\begin{itemize}
    \item \textit{The Adaptations to Strength Training: Morphological and Neurological Contributions to Increased Strength}《力量训练的适应性：增力的形态学与神经学因素》
    \item \textit{A Framework For Understanding the Training Process Leading to Elite Performance}《训练精英运动员的基础框架》
    \item \textit{Exercise Physiology}《训练生理学》
\end{itemize}

\end{document}
